\section*{Coherent-potential repair for a physical restriction tree}

Let $T$ be a finite binary physical restriction tree for a Boolean function
$f$ on the uniform cube.  At an internal node $v$, let $f_v$ be the
restricted Boolean function and let $i(v)$ be a previously unrevealed
coordinate.  Its children are the two physical restrictions
$f_{v,0},f_{v,1}$.  Put $w(v)=2^{-\operatorname{depth}(v)}$.

The B20 notation $D(f_\ell)$ is not defined by the B05--B19 source: there,
$D$ is a packet attached to a chosen selector/data split.  To obtain a
well-posed telescope, first assign a nonnegative scalar potential $P_v$ to
every node, including every terminal node.  Define
\[
  \Gamma_v=P_v-\frac12(P_{v,0}+P_{v,1})
\]
at each internal node.  Then
\[
 \boxed{
 P_{\varnothing}
 =\sum_{v\ \mathrm{internal}}w(v)\Gamma_v
 +\sum_{\ell\ \mathrm{terminal}}w(\ell)P_\ell .}
\tag{1}
\]
This follows by summing the defining recurrences; every non-root,
nonterminal potential occurs once with each sign and the same fair-path
weight.

Total influence obeys the independent physical recursion
\[
 I(f_v)=I_{i(v)}(f_v)+\frac12 I(f_{v,0})+
                         \frac12 I(f_{v,1}).
\]
Consequently,
\[
 \boxed{
 I(f)=\sum_{v\ \mathrm{internal}}w(v)I_{i(v)}(f_v)
 +\sum_{\ell\ \mathrm{terminal}}w(\ell)I(f_\ell).}
\tag{2}
\]
Thus no influence is reused if the tree reveals a coordinate at most once
on each path.  Equations (1)--(2) give the correct sufficient condition:
if, for one coherent preassigned potential,
\[
 \Gamma_v\le C I_{i(v)}(f_v)
 \quad\hbox{and}\quad
 P_\ell\le C I(f_\ell),
\]
then $P_{\varnothing}\le C I(f)$.

For the canonical one-point-defect chain, take $P_{q_n}=D_n$, take the
constant child's potential to be zero, and continue only down the
exceptional child.  Equation (1) then reduces exactly to
\[
 D_N=2^{-(N-2)}D_2+
 \sum_{n=3}^{N}2^{-(N-n)}\Gamma_n^D.
\]
The repair therefore preserves the proved chain identity.  What it does
not supply is a canonical $P_v$ for arbitrary Boolean descendants or the
two displayed bounds.  Those are the first unpaid universal statements.
